% =====================================================================
% PPSN 2026 Poster -- Prior Coverage and Search Bias in Multiview SR
% Inria 2024 charter template (theme InriaChile), DIN A0 portrait.
% =====================================================================
\documentclass[final,svgnames,english]{beamer}
\usepackage[english]{babel}

\usepackage[orientation=portrait,size=a0,scale=1.4]{beamerposter}
\usetheme{InriaPoster}

\usepackage{amsmath,amsthm, amssymb, latexsym}
\usepackage{mathspec}
\setmainfont{Inria Sans}[Scale=MatchLowercase, Kerning=Uppercase]
\setsansfont{Inria Sans}[Scale=MatchLowercase, Kerning=Uppercase]
\setmonofont{Fira Code}[Scale=0.74, Kerning=Uppercase]
%\setmathfont(Digits,Latin,Greek,Symbols)[Scale=MatchLowercase]{Inria Serif}
\setmathrm[Scale=MatchLowercase]{Inria Sans}
%\usepackage{mathastext}

% \boldmath %math en gras

\usepackage{graphicx}
\usepackage{booktabs}
\usepackage{qrcode}
\usepackage{adjustbox}
\usepackage{xspace}
\usepackage{nicefrac}
\usepackage{multiobjective}
\usepackage{fontawesome7}
\usepackage{minted}
\urlstyle{tt}
\usepackage{setspace}

\usepackage{tikz}
\usetikzlibrary{calc, backgrounds, arrows, shapes}
\usepackage{smartdiagram}
\usepackage{lipsum}

\usepackage[backend=biber, style=apa6]{biblatex}
\usepackage{csquotes}

\DeclareFieldFormat{hal}{%
  \mkbibacro{hal}\addcolon\space
  \ifhyperref
    {\href{https://hal.archives-ouvertes.fr/#1}{\nolinkurl{#1}}}
    {\nolinkurl{#1}}}
    
\DeclareFieldAlias{eprint:hal}{hal}
\DeclareFieldAlias{eprint:HAL}{eprint:hal}

\renewbibmacro*{eprint}{%
  \printfield{hal}%
  \newunit\newblock
  \iffieldundef{eprinttype}
    {\printfield{eprint}}
    {\printfield[eprint:\strfield{eprinttype}]{eprint}}}

\DeclareFieldFormat{url}{%
\mkbibacro{doi}\addcolon\space
\ifhyperref
  {\href{https://doi.org/#1}{\nolinkurl{#1}}}
  {\nolinkurl{#1}}}

\DeclareFieldFormat{url}{url\addcolon\space\url{#1}}

\renewbibmacro*{url+urldate}{%
  \ifthenelse{\iffieldundef{url}\OR\NOT\iffieldundef{doi}}
    {}
    {\iffieldundef{urlyear}
      {}
      {%\bibstring{retrieved}%
       \setunit{\addspace}%
       %\printurldate
       %\setunit{\urldatecomma}%
       %\bibstring{from}%
       %\setunit{\addspace}
       }%
     \iffieldundef{url}{}{\printfield{url}\renewcommand*{\finentrypunct}{\relax}}}}

\begin{filecontents*}[noheader,overwrite]{self.bib}
@inproceedings{author-year:keyword,
    title       = {An Amazing Poster Inspired on the Inria Charter},
    author      = {One Author and Another Author},
    booktitle   = {19th {I}nternational {C}onference on {A}mazing {P}osters},
    publisher   = {International Poster Association},
    address     = {Las Condes, Chile},
    year        = {2032},
    doi         = {xx.xxx/xxx-x-xxx-xxxx-x_x},
    hal         = {hal-xyz},
}
\end{filecontents*}
\addbibresource{self.bib} % Bibliography file
%     editor      = {Doina Bucur and Della Cioppa, Antonio  and Ting Hu and Eric Medvet},

\graphicspath{{figures/}}
%\setlength{\parindent}{0pt}
%\renewcommand{\arraystretch}{1.15}

% ---------------------------------------------------------------------
% Highlight boxes. Not specific of the Inria charter but nice to have
% ---------------------------------------------------------------------
\newtcolorbox{callout}{%
  enhanced, sharp corners, colback=inria-2024-rouge!11, width=\linewidth,
  colframe=inria-2024-rouge, boxrule=0pt, leftrule=7.4pt,
  left=11pt, right=5pt, top=9pt, bottom=9pt,
  before skip=11pt, after skip=5pt}

\newtcolorbox{takeaway}{%
  enhanced, sharp corners, colback=inria-2024-vert-tendre!11, width=\linewidth,
  colframe=inria-2024-vert-tendre, boxrule=0pt, leftrule=7.4pt,
  left=11pt, right=5pt, top=9pt, bottom=9pt,
  before skip=11pt, after skip=5pt}
  
\newtcolorbox{inform}{%
  enhanced, sharp corners, colback=inria-2024-gris-bleu!11, width=\linewidth,
  colframe=inria-2024-gris-bleu, boxrule=0pt, leftrule=7.4pt,
  left=11pt, right=5pt, top=9pt, bottom=9pt,
  before skip=11pt, after skip=5pt}

% 
\newtcolorbox{findingblock}[2]{%
  enhanced, sharp corners, colback=inria-2024-#2!1, width=\linewidth,
  colframe=inria-2024-#2, boxrule=0.47pt, leftrule=7.4pt,title={#1}, detach title,
  before upper={\textbf{\tcbtitle}\ },
  coltitle=inria-2024-#2,
  left=11pt, right=5pt, top=9pt, bottom=9pt,
  before skip=5pt, after skip=5pt}


\title{An Amazing Poster Inspired on the Inria Charter}
\author{One Author and Another Author}
\institute[]{Inria Chile Research Center, Las Condes, Santiago, Chile. \url{https://inria.cl} --- \url{first.last@inria.cl}}

\begin{document}
\begin{frame}[fragile]{}
\vspace{-2em}
\begin{columns}[T]
    % Forcing left margin
    \begin{column}{0.01\linewidth}
    \end{column}

    % Content
    \begin{column}{0.98\linewidth}
        % TL;DR section for the main findings. Adjust to your needs!
        \begin{columns}[t]
            \begin{column}{0.32\linewidth}
                \begin{inriablock}{}
                    \large Making nice looking posters (NLP) has been known to be an NP-hard problem.\raggedright
                \end{inriablock}
            \end{column}%\hfill
            \begin{column}{0.32\linewidth}
                \begin{inriablock}{}
                    \large \ldots but, \tcblower \large \textbf{it does not have to be like that!}
                \end{inriablock}
            \end{column}%\hfill
            \begin{column}{0.32\linewidth}
                \begin{inriablock}{}
                    \large We propose a polynomial-time solution to the NLP issue: \textcolor{inria-rouge}{\textbf{contributions are welcomed}}.
                \end{inriablock}
            \end{column}
        \end{columns}
        % Main content
        \begin{columns}[T]
            \begin{column}{.485\linewidth}
            \begin{inriastartblock}{Context}
                \begin{itemize}
                    \item We love \LaTeX\ and we need to make posters.
                    \item This is a poster template based on beamer.
                    \item We have tried to capture the Inria charter with minimal variations.
                    \item Also an starting point, please help and contribute.
                    \item This raises the question: \textbf{what is the meaning of life without posters?}
                \end{itemize}
            \end{inriastartblock}
            
            \begin{inriablock}{Defining the Nice Looking Poster (NLP) Problem}
                Find optimal poster $\hat{p}$ such that
                \[
                  \hat{p} = \argmin_{p\in\set{P}}{\left\lbrace T(p), U(p), F(p)\right\rbrace},\ \text{where}
                \]
                \begin{itemize}
                  \item $T(p)$: time dedicated to prepare the poster,
                  \item $U(p)$: ugliness of the poster,
                  \item $F(p)$: frustration induced by the poster creation process, and
                  \item $\set{P}$: set of all possible posters.
                \end{itemize}
                \begin{takeaway}
                  As we already use \LaTeX\ for our papers, hence we will restrict $\set{P}$ to that domain.
                \end{takeaway}
            \end{inriablock}

            \begin{inriablock}{The building blocks of the Inria poster}
            We have defined four types of blocks:
            \begin{itemize}
                \item \texttt{inriastartblock}: a block to start the poster,
                \item \texttt{inriablock}: blocks to use along the poster,
                \item \texttt{inriafinalblock}: to be used as the final block of the poster, and
                \item \texttt{plainblock}: a simple block for situations where you want to separate content, like in the next block.
            \end{itemize}
            \end{inriablock}

            % a simple block to put context
            \begin{plainblock}{}
            \centering
             %\begin{adjustbox}{max width=0.83\textwidth}
             \input{sample-figs/tikz-sample-figure}
             %\end{adjustbox}
            \end{plainblock}

            \begin{inriablock}{The specific Inria NLP instance}

            In brief, we can say that an Inria NLP is,
            \begin{itemize}
              \item \textbf{Minimalist:} less is more, focus on readability of the content.
              \item \textbf{Sophisticated:} these gradients are so cool.
              \item \textbf{Compliant:} should include the RF Inria logo.
              \begin{itemize}
                \item Refer to the Inria charter of what to do if you want to include some logos.
                \item this second level itemize is cool, right?
              \end{itemize}
            \end{itemize}
            \end{inriablock}

\begin{inriablock}{The (infinite) INLP loop}
%\vspace{-2ex}
\begin{columns}[T]
    \begin{column}{0.45\textwidth}
        \begin{itemize}
              \item Be concise.
              \begin{itemize}
                  \item Avoid long narrative text spans,
                  \item instead itemize or enumerate.
              \end{itemize}
              \item Be precise.
              \item Be illustrative.
              \item Be self-contained.
              \item Ensure readability at a distance.
              \item Stick to the Inria colors.
              \begin{itemize}
                  \item you earn extra points if your plots also use those colors!
              \end{itemize}
            \end{itemize}
    \end{column}
    \begin{column}{0.5\textwidth}
        \centering
        \tikzset{every shadow/.style={fill=none,shadow scale=0}}
\tikzset{module/.append style={shade=none, fill=\col}}
\smartdiagramset{
    border color=\col,
    text color=white,
    module shape=rounded rectangle,
    font=\small,
    text width=0.2\linewidth,
    set color list={inria-2024-framboise,inria-2024-violet,inria-2024-bleu-nuit,inria-2024-bleu-canard,inria-2024-bleu-azur,inria-2024-gris-bleu, inria-2024-cactus,inria-2024-orange},
    circular final arrow disabled=true,
    circular distance=0.38\linewidth,
    arrow tip=to,
    arrow line width=4pt}
\smartdiagram[circular diagram:clockwise]{Paper \texttt{.tex}, Venue rules, Prepare poster, Print, Pin to wall, Present, Take photo, Amaze friends!}
    \end{column}
    \end{columns}
\end{inriablock}

% You can move this block around to fit your poster content
    \begin{plainblock}{}
    \begin{columns}[c]
        \begin{column}{0.092\textwidth}
        \qrcode[height=\textwidth,hyperlink]{https://github.com/Inria-Chile/inria-latex-poster}
        \end{column}
        \begin{column}{0.9\textwidth}%
            {\linespread{0.83}\footnotesize\fullcite{author-year:keyword}\\[-1.28ex]
            \textbf{Code:} \url{https://github.com/Inria-Chile/inria-latex-poster}}
        \end{column}
    \end{columns}
    \end{plainblock}
\end{column}

% Column 2
\begin{column}{.485\linewidth}
    \begin{plainblock}{}
        \begin{findingblock}{Finding 1. Using \texttt{findingblock} to draw attention.}{framboise}
        
        You can use \texttt{findingblock} to draw attention to the main findings of the paper.
        \begin{minted}{latex}
        \begin{findingblock}{Title of finding}{inria-color}
            <text of the finding>
        \end{findingblock}
        \end{minted}
        \begin{itemize}
                \item \textbf{\texttt{inria-color}}: rouge, framboise, violet, bleu-nuit, bleu-canard, bleu-azur, bleu-vert, gris-bleu, cactus, vert-tendre, jaune, orange, and sable.
        \end{itemize}
        \end{findingblock}
        
        \begin{findingblock}{Finding 2. This is actually cool.}{violet}
            \begin{columns}[T]
                \begin{column}{0.6\textwidth}
                    \begin{itemize}
                        \item \lipsum[1][1] 
                        \item \lipsum[1][2]
                        \item \lipsum[1][3]
                        \item \lipsum[1][4] 
                    \end{itemize} 
                \end{column}
                \begin{column}{0.38\textwidth}
                    \begin{tikzpicture}[scale=2.9, every node/.style={font=\small}]
    % Define coordinates for the centers of the 3 circles
    \coordinate (A) at (0, 0.6);
    \coordinate (B) at (-0.7, -0.4);
    \coordinate (C) at (0.7, -0.4);

    % Draw and fill the circles with opacity so intersections show up
    \draw[thick, fill=inria-2024-rouge, fill opacity=0.74] (A) circle (1.2cm);
    \draw[thick, fill=inria-2024-gris-bleu, fill opacity=0.74] (B) circle (1.2cm);
    \draw[thick, fill=inria-2024-framboise, fill opacity=0.74] (C) circle (1.2cm);

    % Outer Set Labels
    \node[color=white] at (0, 1.4) {Inria};
    \node[color=white] at (-1.2, -0.92) {\LaTeX};
    \node[color=white] at (1.2, -0.92) {Poster};

    % 1. Label for the exact center overlap (A intersect B intersect C)
    % Uses barycentric coordinates to find the exact middle of A, B, and C
    \node[text opacity=1, color=white] at (barycentric cs:A=1,B=1,C=1) {You};

    % 2. Labels for dual-overlaps (e.g., just A and B)
    %\node[text opacity=1] at (barycentric cs:A=1,B=1) {$A \cap B$};
    %\node[text opacity=1] at (barycentric cs:A=1,C=1) {$A \cap C$};
    %\node[text opacity=1] at (barycentric cs:B=1,C=1) {$B \cap C$};
\end{tikzpicture}
                \end{column}
            \end{columns}
        \end{findingblock}
    
        \begin{findingblock}{Finding 3. We have other blocks.}{framboise}
            \begin{takeaway}
                \texttt{takeaway}: green block for good news. \lipsum[1][1-2]
            \end{takeaway}
        
            \begin{callout}
                \texttt{callout}: a red block to draw attention. \lipsum[1][1-2]
            \end{callout}
        
            \begin{inform}
                \texttt{inform}: to frame some important information or definition. \lipsum[1][1-2]
            \end{inform}
        \end{findingblock}
    
        \begin{findingblock}{Finding 4. Columns inside your findings are your friend.}{violet}
            \begin{columns}[T]
                \begin{column}{0.38\textwidth}
                    \begin{itemize}
                      \item You can organize your content as (nested) columns.
                            \begin{itemize}
                              \item helps to have narrative next to visual elements, 
                              \item it can be fun.
                            \end{itemize}
                    \end{itemize}
                \end{column}
                \begin{column}{0.6\textwidth}
                    \centering
                    \input{sample-figs/tikz-hepta}
                \end{column}
            \end{columns}
        \end{findingblock}
    \end{plainblock}
    
    \begin{inriafinalblock}{Final remarks}
        \begin{itemize}
            \item Use \verb|\begin{inriafinalblock}{header}...\end{inriafinalblock}| to end your poster.
            \item \textbf{Source for this poster is available in:}
            \begin{itemize}
                \item \url{https://github.com/Inria-Chile/inria-latex-poster}
            \end{itemize} 
            \item \textbf{Inria \LaTeX\ beamer template:}
            \begin{itemize}
                \item \url{https://gitlab.inria.fr/gabarits/latex-beamer-2024}
            \end{itemize}
            \item \textbf{Inria fonts in \LaTeX:}
             \begin{itemize}
                \item \url{https://gitlab.inria.fr/gabarits/latex-inria-fonts}
                \end{itemize}
            \item \textbf{Charter of use of the visual identity of Inria:}
                \begin{itemize}
                    \item \url{https://www.inria.fr/en/charter-use-visual-identity-inria}
                \end{itemize}
            \item \textbf{More info in Numin:}
                \begin{itemize}
                    \item \url{https://numin.inria.fr/portal/g/:spaces:71lwaa/base_de_connaissance/notes/106} 
                \end{itemize}
            \item We are aware that this is far from perfect. 
            \item Feel free to contribute your own improvements.
        \end{itemize}
    \end{inriafinalblock}
    \begin{plainblock}{}
        \centering
        Made with \faHeart\ by Inria Chile -- \url{https://inria.cl}.
    \end{plainblock}
\end{column}
\end{columns}
\end{column}

% right margin
\begin{column}{0.01\linewidth}
\end{column}
\end{columns}
\end{frame}
\end{document}
